% !TEX program = xelatex
\documentclass[lang=zh]{aaron-hw}
\aaronsetup{course=常微分方程, number=3, due=2026/10/28, id=B12345678}
\author{Aaron Lin}

\begin{document}
\maketitle

\begin{problem}[一階線性方程]
解初值問題 $y' + 2y = e^{-x}$，$y(0) = 1$。
\begin{solution}
積分因子為 $\mu(x) = e^{2x}$，所以 $(e^{2x}y)' = e^{x}$，積分得
\[
  y(x) = e^{-x} + C e^{-2x}.
\]
由 $y(0)=1$ 得 $C = 0$，故 $y = e^{-x}$。
\end{solution}
\end{problem}

\begin{problem}
Let $\vb{A} \in \R^{n\times n}$ be symmetric. Show that
\begin{parts}
  \item $\vb{x}\T \vb{A} \vb{x} \le \lambda_{\max} \norm{\vb{x}}^2$ for all $\vb{x}$;
  \item $\argmax_{\norm{\vb{x}} = 1} \vb{x}\T \vb{A} \vb{x}$ is an eigenvector.
\end{parts}
\begin{solution}
Write $\vb{A} = \vb{Q}\vb{\Lambda}\vb{Q}\T$ and set $\vb{y} = \vb{Q}\T\vb{x}$. Then
$\vb{x}\T\vb{A}\vb{x} = \sum_i \lambda_i y_i^2 \le \lambda_{\max}\norm{\vb{y}}^2 = \lambda_{\max}\norm{\vb{x}}^2$.
\begin{proof}
Equality forces $y_i = 0$ whenever $\lambda_i < \lambda_{\max}$.
\end{proof}
\end{solution}
\end{problem}
\end{document}
